\clearpage
\section{消失圈与匹配圈}\label{sec:16}
\subsection{标架}\label{sec:16a}
令 $(M^{2n},\omega)$ 为任意辛流形，$V\subset M$ 为 Lagrangian 球面。$V$ 的\emph{标架}是一个微分同胚 $v:S\to V$，其中 $S$ 为某个 $(n+1)$ 维内积空间的单位球面；$S$ 依赖于 $V$。注意，这与例如 Pontryagin–Thom 构造中带标架流形的概念不同。类似地，若有两个带标架 Lagrangian 球面 $v_k:S_k\to V_k$，$k=0,1$，则它们之间的\emph{带标架同痕} $I=(i,j_0,j_1)$
\SourcePage{229}
由一族 Lagrangian 嵌入 $i^s:S\to M$，$0\leq s\leq1$，以及两个等距映射 $j_k:S\to S_k$ 组成，使 $v_k\circ j_k=i^k$；这里 $S$ 仍是某个内积空间的单位球面。

\StatementNumber{1}\begin{note}\label{st:16.1}
假设 $V_0=V_1=V$ 且 $S_0=S_1=S$ 相同，但有两个不同映射 $v_0,v_1$。显然，若 $v_k$ 在微分同胚空间 $\operatorname{Diff}(S,V)$ 中同痕，则两个带标架 Lagrangian 球面同痕。另一方面，若 $v_k$ 相差 $S$ 的某个反定向等距映射，则令所有 $s$ 的 $i^s=v_0$，并取 $j_0=\Id$、$j_1=v_1^{-1}\circ v_0$，就得到带标架 Lagrangian 球面的同痕。因此，只要关心的仅是同痕类，而我们使用这一概念时基本都是如此，标架信息就归结为 $\pi_0(\operatorname{Diff}(S,V)/O(S))$ 中的一个类。已知 $n<3$ 时 $\operatorname{Diff}(S^n)/O_{n+1}$ 可缩，故这些维数中的标架并不重要。但对较大的 $n$，奇异球面的存在意味着 $\pi_0(\operatorname{Diff}(S^n)/O_{n+1})$ 往往非平凡。\znote{221}{可缩只在 $n<3$ 成立，不要推广成所有维数或写成 $n\leq3$。可缩意味着任意两个标架都同痕，标架信息于是完全消失；$n\geq3$ 时 $\pi_0$ 不再平凡，奇异球面（与 $S^n$ 同胚但不微分同胚的球面）正是其非平凡性的来源。因此后文凡出现“带标架”之处，在 $n\geq3$ 都不能随手丢弃标架，只有在 $n<3$ 才可以。[来源：原书 Remark 16.1，印刷页 221。]}
\end{note}

\subsection{消失圈}\label{sec:16b}
令 $\pi:E^{2n+2}\to S$ 为 Lefschetz 纤维化，仍照例假设正合。\emph{消失路径}是终止于一个临界值的嵌入路径 $\gamma:[0,1]\to S$；更确切地，要求 $\gamma^{-1}(\operatorname{Critv}(\pi))=\{1\}$。对每条这样的路径，可赋予其 \emph{Lefschetz 顶针} $\Delta_\gamma$，即 $E$ 中唯一的嵌入 Lagrangian $(n+1)$ 维球体，满足 $\pi(\Delta_\gamma)=\gamma([0,1])$、$\pi(\partial\Delta_\gamma)=\{\gamma(0)\}$。边界 $V_\gamma=\partial\Delta_\gamma$ 是 $E_{\gamma(0)}$ 中的 Lagrangian 球面，称为 $\gamma$ 的\emph{消失圈}。由于它是 $E$ 中一个 Lagrangian 圆盘的边界，任何消失圈都自动正合；当然，这仅在 $n=1$ 时非平凡。\znote{221}{三点提醒：一、顶针 $\Delta_\gamma$ 是 $(n+1)$ 维球体（实维数），不是球面；其边界 $V_\gamma=\partial\Delta_\gamma$ 才是 $n$ 维球面，二者维数差一。二、$\pi|\Delta_\gamma$ 是到 $\gamma([0,1])$ 的同胚，故 $\pi(\partial\Delta_\gamma)=\{\gamma(0)\}$ 说明整块边界落在起始纤维 $E_{\gamma(0)}$ 内，这正是“消失”一词的几何内容：绕临界值一周，圈便收缩成一点。三、“正合”不是额外假设：由 Stokes 公式，$\theta|_{V_\gamma}=df$，其中 $f$ 是 $\Delta_\gamma$ 上的辛面积函数；$n>1$ 时 $V_\gamma$ 单连通，故并无实质内容。[来源：原书 §(16b)，印刷页 221–222。]}继续之前，记下一个初等事实：

\SourcePage{230}
\StatementNumber{2}\begin{lemma}\label{st:16.2}
对任意 Lagrangian 子流形 $F\subset E$，点 $x\in F$ 是 $\pi$ 的临界点，当且仅当它是 $\pi|F$ 的临界点。
\end{lemma}

对于具有上述性质的圆盘 $\Delta_\gamma$，函数 $p=\gamma^{-1}\circ\pi|\Delta_\gamma:\Delta_\gamma\to[0,1]$ 必在某个内点达到极大值；由\stref{引理}{16.2}，该点必是 $E_{\gamma(1)}$ 中唯一的奇点 $x_1$。类似论证表明，$p$ 在该点的 Hessian 非退化，故 $p$ 是只有一个临界点的 Morse 函数。在 $\Delta_\gamma$ 上赋予一个度量，使其在临界点附近的局部 Morse 坐标中标准化，便可用 $\nabla p$ 的流给出微分同胚
\begin{equation}\label{eq:16.1}
B=B(T_{x_1}\Delta_\gamma)\longrightarrow\Delta_\gamma,
\end{equation}
从 $x_1$ 处切空间的单位球映到整个 Lefschetz 顶针。特别地，限制到边界便得到标架 $v_\gamma:S=\partial B\to V_\gamma$，其在同痕意义下不依赖于度量的选择。将 $v_\gamma$ 视为消失圈结构的一部分，因而消失圈总被视为带标架正合 Lagrangian 球面。

我们从 Lefschetz 顶针的定义性质推出了一些结论，但其存在唯一性仍有待验证。现在来做这件事，先从唯一性开始。以下事实很有用：

\clearpage
\StatementNumber{3}\begin{lemma}\label{st:16.3}
令 $\beta$ 为底空间 $S$ 正则部分中的嵌入路径，$F\subset E$ 为在 $\beta$ 上纤维化的子流形，且每个 $F_{\beta(s)}\subset E_{\beta(s)}$ 都为 Lagrangian。则整个 $F$ 为 Lagrangian，当且仅当沿 $\beta$ 任意一段的平行移动把相应的 $F_{\beta(s)}$ 互相映到彼此。
\end{lemma}

对于 Lefschetz 顶针，这意味着沿 $\gamma$ 的平行移动把水平集 $\Delta_\gamma\cap E_{\gamma(t)}=p^{-1}(t)$，$0\leq t<1$，互相映到彼此。当 $t\to1$ 时，这些水平集收缩到 $\{x_1\}$。因此，$\Delta_\gamma$ 必包含在下述集合中：
\[
\Delta'_\gamma=\{x\in E_{\gamma(t_0)},\ 0\leq t_0<1:
\lim_{t_1\to1}h_{\gamma|[t_0,t_1]}(x)=x_1\}\cup\{x_1\}.
\]
$\Delta'_\gamma$ 的结构未必一眼可见，因为平行移动向量场在 $x_1$ 有极点。为避开这一困难，选择光滑函数 $h:S\to\R$，使 $h(\gamma(t))=0$，且对任意 $t$，$(\gamma'(t),\nabla h_{\gamma(t)})$ 为正定向切向量基。令 $X_H$ 为 $H=h\circ\pi$ 的 Hamilton 向量场。容易看出，它在 $\pi^{-1}(\gamma([0,1)))$ 上与平行移动向量场指向相同方向，同时在 $x_1$ 具有驻点。故 $\Delta'_\gamma$ 包含在
\begin{equation}\label{eq:16.2}
\Delta''_\gamma=\{x\in E_{\gamma(t)},\ 0\leq t\leq1:
\lim_{s\to\infty}\phi_H^s(x)=x_1\}.
\end{equation}
\SourcePage{231}
由于驻点 $x_1$ 双曲，可用稳定流形定理证明，$\Delta''_\gamma$ 是嵌入的 $(n+1)$ 维球体，边界为球面 $\Delta''_\gamma\cap E_{\gamma(0)}$。按假设 $\Delta_\gamma$ 也具有同样性质，所以包含关系 $\Delta_\gamma\subset\Delta'_\gamma\subset\Delta''_\gamma$ 必为等式，证明了 Lefschetz 顶针的唯一性。对于存在性，只须注意\eref{16.2}为 Lagrangian，因为辛流 $\phi_H$ 在其上收缩，且它具有其余全部所需性质。\znote{223}{这是“夹逼”式论证，两处跳步值得点明。其一，$h$ 的选取使 $X_H$ 与平行移动同向，才有 $\Delta'_\gamma\subset\Delta''_\gamma$。其二，$\Delta''_\gamma$ 是嵌入球体、且其边界落在 $E_{\gamma(0)}$ 内，靠的是 $x_1$ 双曲（Hessian 非退化）时稳定流形定理的标准结论，不能只从 $\Delta'_\gamma$ 的定义读出；$x_1$ 处平行移动向量场有极点，正是这一步必须换用 Hamilton 流的原因。存在性反过来取 $\Delta''_\gamma$ 为候选：$\phi_H$ 在其上收缩，故 $\omega$ 沿流的切向消失，从而得到 Lagrangian。[来源：原书 §(16b)，印刷页 222。]}

\StatementNumber{4}\begin{note}\label{st:16.4}
作为上述讨论的副产物，任意一段 $\gamma|[t_0,1]$ 的消失圈与整个路径的消失圈通过平行移动联系：
\begin{equation}\label{eq:16.3}
V_\gamma=h_{\gamma|[0,t_0]}^{-1}(V_{\gamma|[t_0,1]}).
\end{equation}
利用这一观察，可对不必嵌入、但仍满足 $\gamma^{-1}(\operatorname{Critv}(\pi))=\{1\}$ 及 $\gamma'(1)\ne0$ 的路径 $\gamma$ 定义 Lefschetz 顶针和消失圈。具体地，$t_0$ 充分接近 $1$ 时，$\gamma|[t_0,1]$ 为嵌入路径，可用平行移动延伸 $\Delta_{\gamma|[t_0,1]}$。这样定义的 Lefschetz 顶针是浸入而不必嵌入的 Lagrangian 球体；但其边界恰是\eref{16.3}，具有与嵌入情形相同的性质。
\end{note}

\StatementNumber{5}\begin{example}\label{st:16.5}
取\stref{例题}{15.4}定义的简单模型 $\pi:E\to S$。它尚不完全是 Lefschetz 纤维化，但缺少的条件，即全纯凸性，与眼前的问题无关。唯一临界值为 $0$，对任何消失路径 $\gamma$，均可显式确定 Lefschetz 顶针：
\begin{equation}\label{eq:16.4}
\Delta_\gamma=\bigcup_{0\leq t\leq1}\sqrt{\gamma(t)}S^n;
\end{equation}
这里 $\sqrt zS^n=\{x\in\CC^{n+1}:x=\pm\sqrt z\,y\text{，其中某个 }y\in S^n\subset\R^{n+1}\}$。要看出\eref{16.4}成立，可用\stref{例题}{15.4}中的函数 $k$，它在平行移动下不变；再注意 $k^{-1}(0)$ 恰是对所有 $z\in S$ 取子集 $\sqrt zS^n$ 所成的并。
\end{example}
